\section{Eight patterns inside a declared conductor range}\label{sec:patterns}
The conductor sieve proceeded by adjoining correctly normalized columns and retaining both good and bad alternatives at each test prime. The predecessor stages are useful because they show which data each strengthening removed. They should not be read as percentages of Beal solved.

The early twenty five prime depth one system, using inert rational conductors congruent to one modulo seven and bounded by 2000, retained 112 normalized tuples. Further same polynomial tests reduced these to $38,12,6,2$, then a test at $1051$ removed that restricted conductor class. Degree two conductor primes $307,797,1483$, congruent to minus one modulo seven, could not be discarded. Their natural Kummer rows vanish, so they supply genuine extra columns in the necessary system.

The enlarged calculation retained 94 tuples after the tests at $701,911,1009$. Further tests at $1289,1429,1471,1499,1709$ reduced that list to 18. Finally the tests at $2129,2269,2311$ left eight. These are historical exact finite necessary system counts. The fresh reproduction scope and input pins are reported separately in Section~\ref{sec:evidence}. A count quoted from an immutable report is not silently called a new run.

Every final tuple has natural ratio $r=6$, nonzero unit coefficient, and exactly nine ramifying rational primes. Six primes are mandatory:
\[
307,337,547,883,1483,1933.
\]
The other three belong to $197,673,743,797$ in the four support patterns below. Each row has two coefficient conjugate tuples, according to the sign of $u$.
\begin{center}\small
\begin{tabular}{clp{82mm}}\toprule
Pattern&$u$&Ordered beta support and inertia vector\\\midrule
A&$\pm3w$&$307,337,547,673,743,797,883,1483,1933$\newline $(4,4,2,4,3,6,3,1,3)$\\[5pt]
B&$\pm3w$&$197,307,337,547,673,797,883,1483,1933$\newline $(5,1,4,2,6,4,1,4,2)$\\[5pt]
C&$\pm w$&$197,307,337,547,673,743,883,1483,1933$\newline $(1,5,4,1,4,1,6,5,4)$\\[5pt]
D&$\pm w$&$197,307,337,547,743,797,883,1483,1933$\newline $(4,5,5,6,2,4,2,1,6)$\\\bottomrule
\end{tabular}
\end{center}
\figplaceholder{11_support_sets}{Four supports account for the eight coefficient conjugate tuples in the frozen bounded regime. Additional beta ramification changes the column system.}
The table is a necessary cover under the inherited geometry, actual depth one, all inert regime, and the requirement that every rational prime ramifying in beta is at most 2000. It is not a finite cover of all primitive integer triples. Additional divisors of $C$ unramified in beta remain allowed. New beta ramification is outside the cover and must add its columns everywhere, including in the Frobenius evaluations.

\subsection{Why a homogeneous finite sieve can saturate}
Block 70 proved a limitation for the specified zero or good test family. Chebotarev supplies arbitrarily many eligible inert conductor primes where the natural Kummer rows all vanish. On $N$ such primes, the inverse cyclotomic ray character space has dimension $N$ over $\F_7$ under the declared class and unit inputs. Each prescribed test zero value adds at most one linear equation. If $N$ exceeds the number of equations, a nonzero kernel remains. Every trace product at those tests is zero, so the bad alternative accepts the abstract data.

The archive contains an explicit twenty four row witness on twenty five conductor primes above 3000, with a kernel vector having twenty one nonzero inertia coordinates. It is not a global Frey representation. Its purpose is to show that enlarging a fixed list of these homogeneous tests alone cannot give an unrestricted support theorem.

The conclusion is scoped. It does not rule out a finite certificate that uses stronger facts. A necessarily good reduction condition, a nonhomogeneous local value, a justified conductor restriction, or an integral relation from the same actual representation can escape this saturation. The later fixed seven and affine stages seek precisely such information.
